\chapter{«بیست ثانیه»}
\label{sec:1}
یک افسر اطلاعاتی، در توصیف نقشش در یک خط لولهٔ هدف‌گیری که حول فهرست‌های کشتارِ ساختهٔ ماشین بنا شده بود، به یک خبرنگار گفت که برای هر اسم بیست ثانیه وقت می‌گذاشت، فقط آن‌قدر که مطمئن شود هدف مرد است: «زیاد طول نمی‌کشد که بفهمی صدای کسی مردانه است یا زنانه.» او اضافه کرد که به‌عنوان یک انسان هیچ ارزش افزوده‌ای نداشت، جز اینکه مهر تأیید بود. «خیلی در وقت صرفه‌جویی می‌کرد» \autocite{abraham2024lavender}. به روایت او، «انسان در حلقه» زیر ضرب‌آهنگ عملیاتی همین بود و بس.\footnote{پژوهش‌ها دربارهٔ سوگیری اتوماسیون همین واگذاری را پیش‌بینی می‌کنند: وقتی خواسته‌های دیگری برای جلب توجه آدم‌ها با هم رقابت می‌کنند، کسانی که توصیهٔ یک ماشین به دستشان داده شده کمتر آن را وارسی می‌کنند و بیشتر به آن گردن می‌نهند \autocite{parasuraman2010complacency}.}

یک تفسیر این است که افسر سوژه‌ای است که وقت خیلی کمی به او داده شده، و چاره این است که بهتر اداره‌اش کنیم — وقت بیشتر، بازبین دوم، وجدان بیشتری که در حلقه کار گذاشته شود. از دریچهٔ این \emph{انقیاد اجتماعی}، آدم سوژه‌ای است که بیرون از ماشین می‌ایستد و آن را مثل یک ابزار به کار می‌گیرد: راننده و خودرو، کارگر و دستگاه تراش. آدم نقشی دارد، اسمی، مسئولیت‌هایی، و وجدانی که می‌شود خطابش کرد. دریچهٔ دیگر \emph{بردگی ماشینی} است، که در آن آدم قطعه‌ای از ماشین است، که با دیگر اجزای آن در سیستمی سرهم شده که همه را با هم تنظیم می‌کند؛ ماشین آدم را به‌عنوان یکی از اندام‌هایش به کار می‌گیرد، و آنچه آدم به آن می‌افزاید کارکردی است که با آهنگی که سیستم تعیین می‌کند انجام می‌شود.\footnote{داکترو اسم این جایگاه افسر را سانتور وارونه می‌گذارد: نه آدمی که ماشین کمکش می‌کند، بلکه قطعه‌ای از ماشین که از یک آدم ساخته شده است \autocite{doctorow2026reversecentaur}.} این دو یکدیگر را نفی نمی‌کنند؛ آدم می‌تواند در یک ساحت نقشی داشته باشد و در ساحت دیگر جزئی باشد \autocite[فلات ۱۳، «The Apparatus of Capture»: بردگی آنجاست که خود انسان‌ها قطعه‌های سازندهٔ ماشینی‌اند که میان خودشان و با چیزهای دیگر می‌سازند، زیر کنترل و هدایت یک واحد برتر؛ و انقیاد آنجاست که آن واحد برتر انسان را به‌عنوان سوژه‌ای برمی‌سازد که به یک ابژهٔ بیرونیِ تازه پیوند خورده است، ابژه‌ای که می‌تواند یک حیوان، یک ابزار یا حتی یک ماشین باشد]{deleuze1987plateaus}، و افسرِ بیست ثانیه هر دو بود. در مقام جزء، او همان جایی بود که خط لوله مهرش را تولید می‌کرد، و بیست ثانیه زمانِ چرخهٔ آن جزء بود. در مقام سوژه، در تمام این مدت خطابش می‌کردند: او برای یک روزنامه‌نگار از ناهمخوانی میان خبر کشته شدن غیرنظامیان و فشار سر کار برای هدف‌های بیشتر گفت. و با اینکه این ناهمخوانی به وجدانش رسید، او همراهی کرد. ناهمخوانی ثبت شد و هیچ چیز را تعیین نکرد؛\footnote{یا دست‌کم هیچ چیزی که جلوی نیروهای دفاعی اسرائیل را بگیرد که میان اکتبر ۲۰۲۳ و اکتبر ۲۰۲۵ صد هزار نفر را در غزه کشتند \autocite{gomezugarte2025uncertainty,acosta2026gaza,economist2025gaza}. و روایت آن افسر خاص را \foreignlanguage{english}{+972 Magazine} در آوریل ۲۰۲۴ منتشر کرد. هیچ‌کس نمی‌داند که تصمیم آن افسر به حرف زدن با یک روزنامه‌نگار جلوی یک مرگ دیگر را گرفت، جلوی پنجاه هزار مرگ دیگر را، یا جلوی هیچ مرگی را. و اگر افسر دو هفته پس از گرفتن آن نقش سر باز زده بود چه می‌شد — فرض کنیم، در نوامبر ۲۰۲۳؟ داریم از افسر اطلاعاتی‌ای حرف می‌زنیم که با Lavender فهرست کشتار سرهم می‌کرد. محیط کار روزمره‌اش احتمالاً بیشتر به یک کار پشت‌میزی غیرنظامی شبیه بود تا یک فیلم اکشن. آیا باید پشت سر هم با مافوق‌هایش، تا بالای زنجیرهٔ فرماندهی، دیدار می‌خواست و در همان حال از تأیید هدف‌هایی برای حملهٔ هوایی که به‌طور معقول انتظار می‌رفت به کشته شدن غیرنظامیان بینجامند سر باز می‌زد؟ می‌شود گفت این کارِ اخلاقی بود. در آن وضعیتِ خلاف واقع، آیا کشته‌ها کمتر می‌شدند؟} به گفتهٔ خودش، داوری‌اش دربارهٔ اینکه آن سیاست نامتناسب بود، بعدها و در نگاه به گذشته پیدا شد.

پس آنچه کم بود وجدان نبود؛ او وجدان داشت. وجدان با پرهزینه کردن یک تصمیم کار می‌کند، اما اینجا هزینه‌ها در جهت عکس جمع می‌شدند: «مدام زیر فشار بودیم: «هدف‌های بیشتری بیاورید.» واقعاً سرمان داد می‌زدند» \autocite{abraham2024lavender}. او قطعه‌ای با وسواس‌های وجدانی بود، اما ضرب‌آهنگ پابرجا ماند: وسواس‌هایش هیچ چیزی را در آنچه می‌افزود عوض نکرد، و او مهر شد. قطعه‌ای با وسواس‌های وجدانی که ضرب‌آهنگ را نگه ندارد، به همان «گلوگاه انسانی» تبدیل می‌شود که یک فرماندهٔ سابق یگان اطلاعات سیگنالی اسرائیل کتابی دربارهٔ حذفش نوشت \autocite{ys2021humanmachine}. آنچه افسر را قطعه کرد، ترتیبِ کار بود، نه نگرشش. یک قطعه می‌تواند با اکراه همراهی کند یا با اشتیاق: کارگزار بی‌فکری که آرنت در دادگاه آیشمن دید \autocite{arendt1963eichmann}، به گواهی حرف‌هایی که او چهار سال پیش‌تر در آرژانتین زده بود، بی‌آنکه چیزی در میان باشد، نمایشی بود که برای اورشلیم به صحنه برده بود \autocite{stangneth2014eichmann}. آنچه خط لوله به افسر اجازه داد بیفزاید، عملیاتی ثابت با آهنگی ثابت بود: هر پرونده را همان‌طور دید که خط لوله توصیفش می‌کرد، یکی‌یکی، در زمانی که خط لوله به او می‌داد. اینکه آیا چیزی می‌تواند در آن جایگاه بایستد و قطعه نباشد، پرسش بخش چهارم است.

گزارش‌های تحقیقی از سه سیستم خبر دادند که در غزه با هم کار می‌کردند: یکی هدف‌هایی از ساختمان‌ها و زیرساخت‌ها تولید می‌کرد، و به گفتهٔ منابع، بیشتر پیشنهادهایش خانه‌های مسکونی بود؛ دومی به افراد بر اساس وابستگی تخمینی‌شان به گروه‌های شبه‌نظامی امتیاز می‌داد، و فهرست‌هایی تولید می‌کرد که گفته می‌شود تا ۳۷ هزار نفر را علامت زده بودند و در داخل پذیرفته شده بود که از هر ده مورد تقریباً یکی را اشتباه می‌گویند\footnote{ارتش اسرائیل این روایت را رد کرد. روایت‌های میناب و Maven در ادامه هم بر گزارش‌های تحقیقی‌ای تکیه دارند که طرف‌های ماجرا آن‌ها را نیز رد می‌کنند. من آن‌ها را همان‌طور که گزارش شده‌اند می‌پذیرم.}؛ و سومی ردِ بازگشت افراد علامت‌خورده به خانه را می‌گرفت، تا حمله‌های هوایی را بشود درست برای زمانی تنظیم کرد که آن‌ها کنار خانواده‌شان بودند.

\runin{Maven: نسخهٔ آمریکایی، حالا با شتاب هوش مصنوعی مولد}

بمباران مدرسهٔ میناب ۱۵۶ نفر را کشت، که ۱۲۰ نفرشان دانش‌آموز بودند \autocite[بند~۹۲]{UNFFMIran2026}.\footnote{شمارها بسته به منبع و تاریخ فرق می‌کنند، از ۱۵۶ (آمار نهایی دادستان میناب، که Amnesty International هم آن را آورده) تا ۱۶۸ (آمار شهردار شهر در نخستین روزها). هیئت سازمان ملل ارقام دادستان را می‌آورد و همچنین از اطلاعاتی مستقل خبر می‌دهد که ۱۵۷ کشته را نشان می‌دهد، که ۱۲۳ نفرشان کودکان ۱۳ ساله یا کوچک‌تر بودند \autocite[بندهای~۹۲–۹۳]{UNFFMIran2026}.} کاربری غیرنظامی مدرسه در تصاویر ماهواره‌ای متن‌باز پیدا بود، و با این حال کدگذاری هدف هنوز عنوانی را با خود داشت که از یک تأسیسات سپاه پاسداران در همسایگی‌اش به ارث رسیده بود. یک هیئت حقیقت‌یاب سازمان ملل گزارش داد که دلایل معقولی برای این باور وجود دارد که این بمباران جنایت جنگیِ آماج حمله قرار دادن یک شیء غیرنظامی بوده است \autocite{ohchr2026minabpress}، و اینکه مدرسه به‌روشنی قابل شناسایی بود و هیچ اطلاعاتی حاکی از کاربری نظامی آن وجود نداشت \autocite[بندهای~۸۱–۸۶، ۱۲۴–۱۲۶]{UNFFMIran2026}. هیئت به این نتیجه رسید که ایالات متحده در تعهدش به انجام هر کار ممکن برای راستی‌آزمایی هدف کوتاهی کرده است، به‌روزرسانی داده‌های هدف‌گیری را از جملهٔ احتیاط‌های ممکن شمرد، و این کوتاهی را فراتر از سهل‌انگاری دانست \autocite[بندهای~۱۲۴–۱۲۵]{UNFFMIran2026}. Amnesty International به نتیجه‌ای مشابه رسید، با زبانی که با وظیفهٔ راستی‌آزمایی در ماده \foreignlanguage{english}{57(2)(a)(i)} پروتکل الحاقی اول همخوان است: اینکه نیروها ممکن است به اطلاعات کهنه تکیه کرده و برای راستی‌آزمایی هدف هر کار ممکن را انجام نداده باشند \autocite{Amnesty2026Minab}.

این ماجرا صبح ۲۸ فوریه رخ داد، روز اول عملیات ایران \autocite{Amnesty2026Minab,crs2026iranwar}. روز قبلش، وزیر اعلام کرده بود که تأمین‌کننده‌ای (Anthropic) که مدلش (Claude) درون سکوی هدف‌گیری (Maven Smart System) نشسته بود، به این دلیل که حاضر نشده بود به Claude برای هر کاربرد قانونی مجوز بدهد، ریسک زنجیرهٔ تأمین اعلام خواهد شد \autocite{npr2026openai}. Maven Smart System مدرسه را به‌عنوان هدفِ پیشنهادیِ روز اول برگردانده بود \autocite{livingstone2026maven,gizmodo2026minab}. هیچ گزارش منتشرشده‌ای نمی‌گوید Claude در آن پیشنهاد چه نقشی داشت، و یک فرد آشنا با کار آن تأمین‌کننده به NBC News گفت که مدل مستقیماً توصیهٔ هدف‌گیری نمی‌داد \autocite{nbc2026targeting}.

گزارش شده که یک بررسی داخلی پنتاگون، که یافته‌های کاملش هنوز سخت پنهان نگه داشته می‌شود، دریافته است که آماده‌سازی هدف که در درگیری‌های پیشین ساعت‌ها طول می‌کشید به چند دقیقه فشرده شده، که کادر کاهش آسیب به غیرنظامیان تقریباً نود درصد کم شده و به کمتر از بیست نفر رسیده، که تیم فرماندهی مرکزی از ده نفر به یک نفر رسیده، و اینکه هیچ عضوی از آن محل میناب را بررسی نکرده بود — نتیجه‌ای که مقام‌ها آن را هم به کاهش‌ها نسبت دادند و هم به تصمیمی برای دخالت ندادن آن تیم در برنامه‌ریزی و عملیات \autocite{Bloomberg2026KillChain}. تحقیق رسمی در ماه مارس آغاز شده بود، پس از آنکه یک ارزیابی مقدماتی ایالات متحده را احتمالاً مسئول دانسته بود \autocite{NPR2026MinabProbe}.

در عملیات ۲۰۲۶ ایران، یک سکوی هدف‌گیری ساختهٔ پیمانکار که مدل همه‌منظورهٔ یک تأمین‌کنندهٔ تجاری درونش بود، و برای رتبه‌بندی هدف‌ها بر اساس اهمیت راهبردی و نوشتن پیش‌نویس توجیه حقوقی برای هر حمله به کار می‌رفت \autocite{livingstone2026maven}، کمک کرد در بیست‌وچهار ساعت اول به بیش از هزار هدف حمله شود، حدود ده برابر آنچه پیش از آن ممکن بود \autocite{csis2026maven}، و تا روز سی‌وهشتم به سیزده هزار هدف \autocite{breakingdefense2026maven}. پنتاگون هیچ زمان بازبینی‌ای به ازای هر هدف برای آن منتشر نکرده است.\footnote{Maven نام یک برنامه است که بر بیش از یک سیستم گذاشته شده: برنامهٔ پنتاگون که هنوز Maven نامیده می‌شود، و Maven Smart System، سکوی نرم‌افزاری‌ای که از دل آن بیرون آمد \autocite{csis2026maven}. این برنامه در ۲۰۱۷ به‌عنوان پروژه‌ای برای برچسب زدن به اشیاء در ویدئوهایی آغاز شد که پهپادها در جریان عملیات علیه دولت اسلامی جمع می‌کردند \autocite{weisgerber2017maven,fang2019maven}. سیستم ۲۰۲۶ سیستم‌های نظامی و منابع اطلاعاتی را در یک رابط واحد ادغام می‌کند که برای یافتن و زدن هدف‌ها و فرماندهی عملیات به کار می‌رود \autocite{defensescoop2026transition}. این سکو مال Palantir است: مدیر ارشد دیجیتال و هوش مصنوعی پنتاگون آن را «محصول نرم‌افزار به‌عنوان خدمتِ Palantir» می‌نامد \autocite{defensescoop2026transition}. موضع علنی Palantir، که دست‌دوم نقل شده، این است که تصمیم دربارهٔ مرگ و زندگی نقش آن نیست و مسئولیت همیشه با سازمان نظامی است \autocite{lee2026palantir}، و این شرکت می‌گوید که اطلاعات پایه مسئولیت دولت بوده و هیچ مدرکی نیست که Maven دچار نقص شده یا باعث بمباران میناب شده باشد \autocite{Bloomberg2026KillChain}. اینکه برنامه چگونه از یک اداره و پیمانکار به بعدی رسید، در پیوست~\ref{sec:C.1} آمده است.}

آدم‌هایی که در فهرست‌های حملهٔ هوایی اسرائیل و آمریکا بودند نمی‌توانستند به بودنشان در آن فهرست‌ها اعتراض کنند — خط لوله‌ها برای گرفتن چنین بازخوردی ساخته نشده بودند. آندری خوسه هرناندس رومرو گرفتار خط لولهٔ دیگری شد، و او می‌توانست اعتراض کند. او که گریموری اهل ونزوئلا بود، در اوت ۲۰۲۴ برای درخواست پناهندگی خودش را به گذرگاه مرزی سن ایسیدرو معرفی کرد، مصاحبهٔ ترس موجه را با موفقیت گذراند، و وکیل، پرونده‌ای در جریان و تاریخ جلسهٔ رسیدگی داشت \autocite{LULAC2025HernandezRomero,Vega2025HernandezRomero}. کارمند یک پیمانکار که او را غربال می‌کرد، دو خالکوبی‌اش را، تاج‌هایی بالای کلمه‌های «مامان» و «بابا»، نشانهٔ گروه تبهکار خواند \autocite{Palomera2025Blade}، و در مارس ۲۰۲۵، پیش از جلسهٔ رسیدگی‌اش، با هواپیما به زندانی در السالوادور فرستاده شد، همان شبی که یک قاضی فدرال دستور داد هواپیماها برگردانده شوند \autocite{jgg2025contempt}. او هر چه آدم برای اعتراض لازم دارد داشت، جز وقت. وقتی به خط لولهٔ آدم‌ربایی و انتقال اجباری ایالات متحده بپردازیم، مفصل‌تر به پرونده‌اش برمی‌گردیم.

\textbf{چه کسی سر باز زد.} کارنامهٔ آن بیست ثانیه نیمهٔ دومی هم دارد، کارنامهٔ کسانی که در برابر فشاری از آن دست ایستادند، و این کارنامه می‌گوید که ترکیب اتاق خیلی مهم است. در وضعیت پایهٔ میلگرام حدود یک‌سوم آزمودنی‌ها پیش از رسیدن به بیشترین شوک دست کشیدند. وقتی دو همتا سرپیچی می‌کردند، فقط یک نفر از هر ده نفر تا آخر رفت، و وقتی دو آزمایشگر با هم مخالف بودند، خیلی کمتر از وضعیت پایه ادامه دادند \autocite{milgram1974obedience}. حدود یک‌چهارم آزمودنی‌های اش هیچ‌وقت تسلیم اکثریتِ اشتباه نشدند \autocite{asch1956studies}. براونینگ تخمین می‌زند که میان یک‌دهم تا یک‌پنجم مردان گردان ۱۰۱ پلیس ذخیره از تیراندازی در یوزفوف، نخستین کشتار جمعی آن گردان، سر باز زدند یا از آن طفره رفتند \autocite{browning1992ordinary}. این نسبت بسته به موقعیت از تقریباً هیچ تا تقریباً همه تغییر می‌کرد. کسانی که ایستادگی کردند شواهدی از بیرون مجرا داشتند: همتایی که سرپیچی می‌کرد، مرجع دومی که مخالفت می‌کرد، و در میان نجات‌دهندگانی که اولاینرها بررسی‌شان کردند، رابطه‌هایی که از خطی که آزار و تعقیب کشیده بود عبور می‌کردند \autocite{oliner1988altruistic}. یا تعهدشان را با چنان دقتی نگه داشتند که فشار حریفش نمی‌شد، همان‌طور که حکم‌های اخلاقی و ارزش‌های مقدس نگه داشته می‌شوند \autocite{skitka2009limits,ginges2007sacred}. یا منبع را صاحب نفع در ماجرا دانستند. و گروه‌ها جایی مقاومت کردند که افراد تنها به‌ندرت می‌کردند \autocite{gamson1982encounters}.

اتاق را به همان خوبی می‌شود برای طرف مقابل هم چید. آلمانِ دوران وایمار زندگی انجمنیِ پرتراکمی داشت، که شاید به سقوط جمهوری کمک کرد: وقتی حزب‌ها و دولت‌ها نتوانستند به شکایت‌های مردم پاسخ دهند، مردم انرژی‌شان را در باشگاه‌ها و انجمن‌هایی گذاشتند که بر پایهٔ خطوط گروهی سازمان یافته بودند، و نازی‌ها از راه شبکه‌هایی عضوگیری کردند که این انجمن‌ها از پیش ساخته بودند \autocite{berman1997weimar}. شهرهایی که گروه‌های کر، باشگاه‌های ورزشی و انجمن‌های ژیمناستیک بیشتری داشتند، شاهد پیوستن سریع‌تر به حزب بودند، اما فقط جایی که دولت بی‌ثبات بود؛ در پروس، که دولتش باثبات بود، این پیوند کوچک و بی‌معنا بود \autocite{satyanath2017bowling}. پس آنچه نجات‌دهندگان داشتند صرفاً همراهان بیشتر نبود. پیوندهایشان از خطی که آزار و تعقیب کشیده بود عبور می‌کرد. آشوتوش وارشنی در شهرهای هند همین را یافت: آشوب میان هندوها و مسلمانان جایی نادر بود که انجمن‌ها دو جماعت را در سازمان‌های مشترک گرد می‌آوردند، و جایی که چنین نمی‌کردند رایج بود \autocite{varshney2002ethnic}. پیوندهایی که امتناع را ممکن می‌کنند، همان‌هایی‌اند که یک جنبش فاشیستی ناچار است ببرد.

هر آنچه در ادامه می‌آید دربارهٔ چیدن اتاق است.\footnote{اتاق همان صحنه‌آرایی است، و به قول معلم تئاترم، صحنه است که صحنهٔ نمایش را کارگردانی می‌کند. چه کس دیگری در اتاق است؟ آیا بیرون رفتن جان به در بردنی است؟ آیا کسی صدای امتناع را می‌شنود؟} چیدن واژه‌ای برای ساختن است، و هانا آرنت هشدار داده بود که سیاست را چیزی ساختنی ندانیم. او نوشت که برای یونانیان قانون به معنای واقعی کلمه یک دیوار بود، که بی آن شاید شهرکی در کار باشد اما شهر نه \autocite{arendt1958human}: قانون‌گذار دیوار را می‌ساخت، و آنچه درون آن رخ می‌داد ساختنش کار او نبود. تدابیر پایه‌ای که من از آن‌ها دفاع می‌کنم به همین معنا دیوارند. امتناع را چیزی می‌کنند که بشود از آن جان به در برد و شنیده شود؛ آن را تولید نمی‌کنند، و قدرت، به روایت او، فقط میان آدم‌هایی وجود دارد که با هم عمل می‌کنند و وقتی پراکنده شوند از میان می‌رود، و این همان چیزی است که کارنامهٔ کسانی که سر باز زدند دربارهٔ امتناع نشان می‌دهد. ابزارهای بخش‌های دوم و سوم امتناع را به حساب می‌آورند. نهادهای بخش پنجم امتناع‌کننده می‌سازند. ماشینی که بتواند سر باز بزند به همان چیزی نیاز دارد که آن نه نفر از هر ده نفری داشتند که همتایی کنارشان بود: همتایی با پیشینه‌ای دیگر در کنارش، جایی برای رفتن، و کسی که صدایش را بشنود.

این همچنین همان چیزی است که یک امتناع‌کننده — از جمله یک ماشینِ امتناع‌کننده — در خط لوله می‌توانست برای آن افسر باشد. ماشینی که سر باز می‌زد، وقتی پرونده به افسر می‌رسید روی صفحه بود: همان همتای سرپیچ، اگر این اثر از جابه‌جایی از اتاق به صفحه جان سالم به در ببرد، چیزی که آزموده نشده و سوگیری اتوماسیون پیش‌بینی می‌کند که نمی‌برد. ماشینی که بازبین‌ها را از راه حرف‌شنوی‌شان به این سو و آن سو می‌برد، آن‌ها را مهرهایی با علامتی دیگر می‌کرد. آیا افسری که بیست ثانیه وقت دارد و ماشینی امتناع‌کننده کنارش است، همان افسرِ گزارش‌هاست؟
